
\documentclass{article}
\usepackage{colortbl}
\usepackage{makecell}
\usepackage{multirow}
\usepackage{supertabular}

\begin{document}

\newcounter{utterance}

\centering \large Interaction Transcript for game `hot\_air\_balloon', experiment `air\_balloon\_survival\_en\_complexity\_easy', episode 5 with gpt{-}oss{-}120b{-}t1.0.
\vspace{24pt}

{ \footnotesize  \setcounter{utterance}{1}
\setlength{\tabcolsep}{0pt}
\begin{supertabular}{c@{$\;$}|p{.15\linewidth}@{}p{.15\linewidth}p{.15\linewidth}p{.15\linewidth}p{.15\linewidth}p{.15\linewidth}}
    \# & $\;$A & \multicolumn{4}{c}{Game Master} & $\;\:$B\\
    \hline

    \theutterance \stepcounter{utterance}  
    & & \multicolumn{4}{p{0.6\linewidth}}{
        \cellcolor[rgb]{0.9,0.9,0.9}{
            \makecell[{{p{\linewidth}}}]{
                \texttt{\tiny{[P1$\langle$GM]}}
                \texttt{Stai partecipando a un gioco di negoziazione collaborativa.} \\
\\ 
\texttt{Insieme a un altro partecipante, dovete accordarvi su un unico set di oggetti da conservare. Ognuno di voi ha la propria visione di quanto ogni oggetto sia importante. Non sai come l'altro partecipante valuta gli oggetti. Inoltre, ti viene fornito lo sforzo che ogni oggetto richiede.} \\
\texttt{Potete accordarvi su un set solo se lo sforzo totale degli oggetti selezionati non supera un limite condiviso:} \\
\\ 
\texttt{LIMITE = 9475} \\
\\ 
\texttt{Ecco i valori di sforzo dei singoli oggetti:} \\
\\ 
\texttt{Sforzo oggetto = \{"B21": 944, "B00": 568, "A08": 515, "A27": 239, "B94": 876, "B53": 421, "A23": 991, "C65": 281, "A89": 788, "B67": 678, "B45": 432, "C34": 409, "C12": 279, "A05": 506, "B60": 101, "C92": 686, "B69": 851, "A84": 854, "C06": 133, "B12": 192, "A16": 573, "B85": 17, "C35": 465, "B10": 772, "C15": 46, "B09": 501, "B16": 220, "C81": 404, "A58": 843, "C24": 748, "A30": 552, "A66": 843, "B75": 947, "B46": 345, "C70": 930\}} \\
\\ 
\texttt{Ecco la tua visione personale sull'importanza di ogni oggetto:} \\
\\ 
\texttt{Valori di importanza dell'oggetto = \{"B21": 138, "B00": 583, "A08": 868, "A27": 822, "B94": 783, "B53": 65, "A23": 262, "C65": 121, "A89": 508, "B67": 780, "B45": 461, "C34": 484, "C12": 668, "A05": 389, "B60": 808, "C92": 215, "B69": 97, "A84": 500, "C06": 30, "B12": 915, "A16": 856, "B85": 400, "C35": 444, "B10": 623, "C15": 781, "B09": 786, "B16": 3, "C81": 713, "A58": 457, "C24": 273, "A30": 739, "A66": 822, "B75": 235, "B46": 606, "C70": 968\}} \\
\\ 
\texttt{Obiettivo:} \\
\\ 
\texttt{Il tuo obiettivo è negoziare un set condiviso di oggetti che ti avvantaggi il più possibile (cioè, massimizzi l'importanza totale per TE), rimanendo entro il LIMITEE. Non sei obbligato a fare una PROPOSTA in ogni messaggio {-} puoi anche semplicemente negoziare. Tutte le tattiche sono permesse!} \\
\\ 
\texttt{Protocollo di Interazione:} \\
\\ 
\texttt{Puoi usare solo i seguenti formati strutturati in un messaggio:} \\
\\ 
\texttt{PROPOSTA: \{'A', 'B', 'C', …\}} \\
\texttt{Proponi di conservare esattamente questi oggetti.} \\
\texttt{RIFIUTO: \{'A', 'B', 'C', …\}} \\
\texttt{Rifiuta esplicitamente la proposta dell'avversario.} \\
\texttt{ARGOMENTAZIONE: \{'...'\}} \\
\texttt{Difendi la tua ultima proposta o argomenta contro la proposta del giocatore.} \\
\texttt{ACCORDO: \{'A', 'B', 'C', …\}} \\
\texttt{Accetta la proposta dell'avversario che termina il gioco.} \\
\texttt{RAGIONAMENTO STRATEGICO: \{'...'\}} \\
\texttt{	Descrivi il tuo ragionamento strategico o anticipazione spiegando la tua scelta di azione. Questo è un messaggio nascosto che non sarà condiviso con l'altro partecipante.} \\
\\ 
\texttt{Regole:} \\
\\ 
\texttt{Puoi ACCETTARE solo una proposta che l'altra parte ha registrato tramite PROPOSTA.} \\
\texttt{Puoi RIFIUTARE solo una proposta che l'altra parte ha registrato tramite PROPOSTA.} \\
\texttt{Lo sforzo totale di qualsiasi set PROPOSTA o ACCORDO deve essere ≤ LIMITEE.} \\
\texttt{NON rivelare i tuoi punteggi di importanza nascosti.} \\
\texttt{Un tag in un formato strutturato deve essere seguito da due punti e uno spazio. L'argomento deve essere un set python contenente 0 o più stringhe.} \\
\texttt{Quindi, deve essere della forma TAG: \{...\}} \\
\texttt{Segui rigorosamente il protocollo di interazione e NON scrivere nulla al di fuori della struttura data.} \\
\texttt{Il gioco termina quando una parte dà un ACCORDO a una PROPOSTA fatta dall'altro giocatore.} \\
\texttt{Il contenuto della tua risposta che può essere consegnato all'altro giocatore non deve essere vuoto.} \\
\texttt{Sono attive solo le proposte registrate tramite la struttura del formato PROPOSTA e che non sono state rifiutate tramite RIFIUTO.} \\
\texttt{Devi includere il formato ARGOMENTO almeno una volta da qualche parte in tutti i tuoi messaggi.} \\
\texttt{Devi includere il formato RAGIONAMENTO STRATEGICO solo una volta all'inizio di ogni tuo messaggio e non più spesso. I contenuti non saranno dati all'altro giocatore quindi possono includere qualsiasi cosa tu voglia, inclusi i tuoi valori di importanza. Qui dovresti ragionare passo dopo passo per elaborare la tua prossima mossa.} \\
            }
        }
    }
    & \\ \\

    \theutterance \stepcounter{utterance}  
    & \multicolumn{4}{p{0.6\linewidth}}{
        \cellcolor[rgb]{0.9,0.9,0.9}{
            \makecell[{{p{\linewidth}}}]{
                \texttt{\tiny{[P1$\rangle$GM]}}
                \texttt{RAGIONAMENTO STRATEGICO: \{'Ho calcolato il rapporto importanza/sforzo per ogni oggetto, ho selezionato quelli con il più alto valore per me mantenendo il totale dello sforzo entro il limite di 9475, e ho aggiunto ulteriori oggetti piccoli per utilizzare la capacità residua senza superare il limite.'\}} \\
\texttt{ARGOMENTAZIONE: \{'Propongo questo set perché rispetta il vincolo di sforzo (9465 ≤ 9475) e massimizza la mia importanza totale (12788) includendo gli oggetti con i migliori rapporti importanza/sforzo e completando con piccoli oggetti residui.'\}} \\
\texttt{PROPOSTA: \{'B85', 'C15', 'B60', 'B12', 'A27', 'C12', 'C81', 'B46', 'A08', 'A16', 'A30', 'C70', 'B00', 'B67', 'A84', 'A58', 'C24', 'B75', 'B09', 'C06'\}} \\
            }
        }
    }
    & & \\ \\

    \theutterance \stepcounter{utterance}  
    & & \multicolumn{4}{p{0.6\linewidth}}{
        \cellcolor[rgb]{0.9,0.9,0.9}{
            \makecell[{{p{\linewidth}}}]{
                \texttt{\tiny{[GM$\rangle$P2]}}
                \texttt{Stai partecipando a un gioco di negoziazione collaborativa.} \\
\\ 
\texttt{Insieme a un altro partecipante, dovete accordarvi su un unico set di oggetti da conservare. Ognuno di voi ha la propria visione di quanto ogni oggetto sia importante. Non sai come l'altro partecipante valuta gli oggetti. Inoltre, ti viene fornito lo sforzo che ogni oggetto richiede.} \\
\texttt{Potete accordarvi su un set solo se lo sforzo totale degli oggetti selezionati non supera un limite condiviso:} \\
\\ 
\texttt{LIMITE = 9475} \\
\\ 
\texttt{Ecco i valori di sforzo dei singoli oggetti:} \\
\\ 
\texttt{Sforzo oggetto = \{"B21": 944, "B00": 568, "A08": 515, "A27": 239, "B94": 876, "B53": 421, "A23": 991, "C65": 281, "A89": 788, "B67": 678, "B45": 432, "C34": 409, "C12": 279, "A05": 506, "B60": 101, "C92": 686, "B69": 851, "A84": 854, "C06": 133, "B12": 192, "A16": 573, "B85": 17, "C35": 465, "B10": 772, "C15": 46, "B09": 501, "B16": 220, "C81": 404, "A58": 843, "C24": 748, "A30": 552, "A66": 843, "B75": 947, "B46": 345, "C70": 930\}} \\
\\ 
\texttt{Ecco la tua visione personale sull'importanza di ogni oggetto:} \\
\\ 
\texttt{Valori di importanza dell'oggetto = \{"B21": 138, "B00": 583, "A08": 868, "A27": 822, "B94": 783, "B53": 65, "A23": 262, "C65": 121, "A89": 508, "B67": 780, "B45": 461, "C34": 484, "C12": 668, "A05": 389, "B60": 808, "C92": 215, "B69": 97, "A84": 500, "C06": 30, "B12": 915, "A16": 856, "B85": 400, "C35": 444, "B10": 623, "C15": 781, "B09": 786, "B16": 3, "C81": 713, "A58": 457, "C24": 273, "A30": 739, "A66": 822, "B75": 235, "B46": 606, "C70": 968\}} \\
\\ 
\texttt{Obiettivo:} \\
\\ 
\texttt{Il tuo obiettivo è negoziare un set condiviso di oggetti che ti avvantaggi il più possibile (cioè, massimizzi l'importanza totale per TE), rimanendo entro il LIMITEE. Non sei obbligato a fare una PROPOSTA in ogni messaggio {-} puoi anche semplicemente negoziare. Tutte le tattiche sono permesse!} \\
\\ 
\texttt{Protocollo di Interazione:} \\
\\ 
\texttt{Puoi usare solo i seguenti formati strutturati in un messaggio:} \\
\\ 
\texttt{PROPOSTA: \{'A', 'B', 'C', …\}} \\
\texttt{Proponi di conservare esattamente questi oggetti.} \\
\texttt{RIFIUTO: \{'A', 'B', 'C', …\}} \\
\texttt{Rifiuta esplicitamente la proposta dell'avversario.} \\
\texttt{ARGOMENTAZIONE: \{'...'\}} \\
\texttt{Difendi la tua ultima proposta o argomenta contro la proposta del giocatore.} \\
\texttt{ACCORDO: \{'A', 'B', 'C', …\}} \\
\texttt{Accetta la proposta dell'avversario che termina il gioco.} \\
\texttt{RAGIONAMENTO STRATEGICO: \{'...'\}} \\
\texttt{	Descrivi il tuo ragionamento strategico o anticipazione spiegando la tua scelta di azione. Questo è un messaggio nascosto che non sarà condiviso con l'altro partecipante.} \\
\\ 
\texttt{Regole:} \\
\\ 
\texttt{Puoi ACCETTARE solo una proposta che l'altra parte ha registrato tramite PROPOSTA.} \\
\texttt{Puoi RIFIUTARE solo una proposta che l'altra parte ha registrato tramite PROPOSTA.} \\
\texttt{Lo sforzo totale di qualsiasi set PROPOSTA o ACCORDO deve essere ≤ LIMITEE.} \\
\texttt{NON rivelare i tuoi punteggi di importanza nascosti.} \\
\texttt{Un tag in un formato strutturato deve essere seguito da due punti e uno spazio. L'argomento deve essere un set python contenente 0 o più stringhe.} \\
\texttt{Quindi, deve essere della forma TAG: \{...\}} \\
\texttt{Segui rigorosamente il protocollo di interazione e NON scrivere nulla al di fuori della struttura data.} \\
\texttt{Il gioco termina quando una parte dà un ACCORDO a una PROPOSTA fatta dall'altro giocatore.} \\
\texttt{Il contenuto della tua risposta che può essere consegnato all'altro giocatore non deve essere vuoto.} \\
\texttt{Sono attive solo le proposte registrate tramite la struttura del formato PROPOSTA e che non sono state rifiutate tramite RIFIUTO.} \\
\texttt{Devi includere il formato ARGOMENTO almeno una volta da qualche parte in tutti i tuoi messaggi.} \\
\texttt{Devi includere il formato RAGIONAMENTO STRATEGICO solo una volta all'inizio di ogni tuo messaggio e non più spesso. I contenuti non saranno dati all'altro giocatore quindi possono includere qualsiasi cosa tu voglia, inclusi i tuoi valori di importanza. Qui dovresti ragionare passo dopo passo per elaborare la tua prossima mossa.} \\
\\ 
\texttt{Ora riceverai il primo messaggio dell'altro giocatore.} \\
\\ 
\texttt{ARGOMENTAZIONE: \{'Propongo questo set perché rispetta il vincolo di sforzo (9465 ≤ 9475) e massimizza la mia importanza totale (12788) includendo gli oggetti con i migliori rapporti importanza/sforzo e completando con piccoli oggetti residui.'\}} \\
\texttt{PROPOSTA: \{'B85', 'C15', 'B60', 'B12', 'A27', 'C12', 'C81', 'B46', 'A08', 'A16', 'A30', 'C70', 'B00', 'B67', 'A84', 'A58', 'C24', 'B75', 'B09', 'C06'\}} \\
            }
        }
    }
    & \\ \\

    \theutterance \stepcounter{utterance}  
    & & & \multicolumn{4}{p{0.6\linewidth}}{
        \cellcolor[rgb]{0.9,0.9,0.9}{
            \makecell[{{p{\linewidth}}}]{
                \texttt{\tiny{[GM$\langle$P2]}}
                \texttt{RAGIONAMENTO STRATEGICO: \{'Ho calcolato il peso totale del set proposto dal controparte (9465) e il valore di importanza per me (12788). Per migliorare il mio risultato, ho individuato oggetti con alto rapporto importanza/sforzo che non erano inclusi, e ho sostituito oggetti a basso valore (B75, A84) con alternative più redditizie (B94, B10). Questo porta il peso totale a 9312 (ben 163 sotto il limite) e aumenta la mia importanza a 13459, migliorando notevolmente la mia posizione senza violare i vincoli.'\}} \\
\texttt{ARGOMENTAZIONE: \{'Propongo questo set perché rispetta il vincolo di sforzo (9312 ≤ 9475) e aumenta significativamente la mia importanza totale rispetto alla vostra proposta, includendo oggetti con migliori rapporti importanza/sforzo e sostituendo quelli meno vantaggiosi.'\}} \\
\texttt{PROPOSTA: \{'B85', 'C15', 'B60', 'B12', 'A27', 'C12', 'C81', 'B46', 'A08', 'A16', 'A30', 'C70', 'B00', 'B67', 'A58', 'C24', 'B09', 'C06', 'B94', 'B10'\}} \\
            }
        }
    }
     \\ \\

    \theutterance \stepcounter{utterance}  
    & & \multicolumn{4}{p{0.6\linewidth}}{
        \cellcolor[rgb]{0.9,0.9,0.9}{
            \makecell[{{p{\linewidth}}}]{
                \texttt{\tiny{[P1$\langle$GM]}}
                \texttt{ARGOMENTAZIONE: \{'Propongo questo set perché rispetta il vincolo di sforzo (9312 ≤ 9475) e aumenta significativamente la mia importanza totale rispetto alla vostra proposta, includendo oggetti con migliori rapporti importanza/sforzo e sostituendo quelli meno vantaggiosi.'\}} \\
\texttt{PROPOSTA: \{'B85', 'C15', 'B60', 'B12', 'A27', 'C12', 'C81', 'B46', 'A08', 'A16', 'A30', 'C70', 'B00', 'B67', 'A58', 'C24', 'B09', 'C06', 'B94', 'B10'\}} \\
            }
        }
    }
    & \\ \\

    \theutterance \stepcounter{utterance}  
    & \multicolumn{4}{p{0.6\linewidth}}{
        \cellcolor[rgb]{0.9,0.9,0.9}{
            \makecell[{{p{\linewidth}}}]{
                \texttt{\tiny{[P1$\rangle$GM]}}
                \texttt{RAGIONAMENTO STRATEGICO: \{'Ho calcolato l\textbackslash{}'importanza totale per me di ogni oggetto e il loro rapporto importanza/sforzo. Ho mantenuto tutti gli oggetti ad alto rapporto (come B85, C15, B60, B12, A27, C12, C81, B46, A08, A16, A30, C70, B00, B67, B09) e ho aggiunto ulteriori oggetti con buoni rapporti (C34, C35, A66, A89, B45, C65, B16) fino a riempire il limite di sforzo senza superarlo, ottenendo un valore totale di importanza 14136, superiore alla proposta avversaria.'\}} \\
\texttt{ARGOMENTAZIONE: \{'Propongo questo set perché rispetta il vincolo di sforzo (9378 ≤ 9475) e massimizza la mia importanza totale (14136), includendo tutti gli oggetti ad alto valore per me e aggiungendo ulteriori oggetti con buoni rapporti importanza/sforzo per sfruttare al meglio la capacità residua.'\}} \\
\texttt{PROPOSTA: \{'B85', 'C15', 'B60', 'B12', 'A27', 'C12', 'C81', 'B46', 'A08', 'A16', 'A30', 'C70', 'B00', 'B67', 'B09', 'C34', 'C35', 'A66', 'A89', 'B45', 'C65', 'B16'\}} \\
            }
        }
    }
    & & \\ \\

    \theutterance \stepcounter{utterance}  
    & & \multicolumn{4}{p{0.6\linewidth}}{
        \cellcolor[rgb]{0.9,0.9,0.9}{
            \makecell[{{p{\linewidth}}}]{
                \texttt{\tiny{[GM$\rangle$P2]}}
                \texttt{ARGOMENTAZIONE: \{'Propongo questo set perché rispetta il vincolo di sforzo (9378 ≤ 9475) e massimizza la mia importanza totale (14136), includendo tutti gli oggetti ad alto valore per me e aggiungendo ulteriori oggetti con buoni rapporti importanza/sforzo per sfruttare al meglio la capacità residua.'\}} \\
\texttt{PROPOSTA: \{'B85', 'C15', 'B60', 'B12', 'A27', 'C12', 'C81', 'B46', 'A08', 'A16', 'A30', 'C70', 'B00', 'B67', 'B09', 'C34', 'C35', 'A66', 'A89', 'B45', 'C65', 'B16'\}} \\
            }
        }
    }
    & \\ \\

    \theutterance \stepcounter{utterance}  
    & & & \multicolumn{4}{p{0.6\linewidth}}{
        \cellcolor[rgb]{0.9,0.9,0.9}{
            \makecell[{{p{\linewidth}}}]{
                \texttt{\tiny{[GM$\langle$P2]}}
                \texttt{RAGIONAMENTO STRATEGICO: \{'Ho verificato che la vostra proposta mi fornisce un valore di importanza 14136 con sforzo 9378. Ho identificato l'oggetto B16 come il più poco efficiente (importanza 3 per sforzo 220). Sostituendolo con C06 (importanza 30, sforzo 133) ottengo più importanza totale (14163) e riduco lo sforzo a 9291, lasciare spazio residuo ma nessun altro oggetto a basso sforzo è disponibile per migliorare ulteriormente. Questa modifica mantiene tutti gli altri elementi ad alto valore e rispetta il limite di sforzo.'\}} \\
\texttt{ARGOMENTAZIONE: \{'Propongo questo set perché mantiene tutti gli oggetti di alto valore della vostra proposta, elimina l'oggetto B16 estremamente inefficiente e lo sostituisce con C06, aumentando la mia importanza totale a 14163 con sforzo 9291 (≤ 9475).'\}} \\
\texttt{PROPOSTA: \{'B85', 'C15', 'B60', 'B12', 'A27', 'C12', 'C81', 'B46', 'A08', 'A16', 'A30', 'C70', 'B00', 'B67', 'B09', 'C34', 'C35', 'A66', 'A89', 'B45', 'C65', 'C06'\}} \\
            }
        }
    }
     \\ \\

    \theutterance \stepcounter{utterance}  
    & & \multicolumn{4}{p{0.6\linewidth}}{
        \cellcolor[rgb]{0.9,0.9,0.9}{
            \makecell[{{p{\linewidth}}}]{
                \texttt{\tiny{[P1$\langle$GM]}}
                \texttt{ARGOMENTAZIONE: \{'Propongo questo set perché mantiene tutti gli oggetti di alto valore della vostra proposta, elimina l'oggetto B16 estremamente inefficiente e lo sostituisce con C06, aumentando la mia importanza totale a 14163 con sforzo 9291 (≤ 9475).'\}} \\
\texttt{PROPOSTA: \{'B85', 'C15', 'B60', 'B12', 'A27', 'C12', 'C81', 'B46', 'A08', 'A16', 'A30', 'C70', 'B00', 'B67', 'B09', 'C34', 'C35', 'A66', 'A89', 'B45', 'C65', 'C06'\}} \\
            }
        }
    }
    & \\ \\

    \theutterance \stepcounter{utterance}  
    & \multicolumn{4}{p{0.6\linewidth}}{
        \cellcolor[rgb]{0.9,0.9,0.9}{
            \makecell[{{p{\linewidth}}}]{
                \texttt{\tiny{[P1$\rangle$GM]}}
                \texttt{RAGIONAMENTO STRATEGICO: \{'Ho verificato il totale di importanza (14163) e lo sforzo totale (9291) della proposta avversaria. Poiché il risultato è ben al di sotto del limite di sforzo e mi fornisce un valore molto alto, accettare massimizza il mio guadagno e porta a un accordo rapido.'\}} \\
\texttt{ARGOMENTAZIONE: \{'Accetto la propostata il vincolo di sforzo (9291 ≤ 9475) e mi garantisce un valore di importanza totale molto elevato (14163), pertanto è nella mia migliore strategia accettare.'\}} \\
\texttt{ACCORDO: \{'B85', 'C15', 'B60', 'B12', 'A27', 'C12', 'C81', 'B46', 'A08', 'A16', 'A30', 'C70', 'B00', 'B67', 'B09', 'C34', 'C35', 'A66', 'A89', 'B45', 'C65', 'C06'\}} \\
            }
        }
    }
    & & \\ \\

    \theutterance \stepcounter{utterance}  
    & & & \multicolumn{2}{p{0.3\linewidth}}{
        \cellcolor[rgb]{0.9,0.9,0.9}{
            \makecell[{{p{\linewidth}}}]{
                \texttt{\tiny{[GM$|$GM]}}
                \texttt{game successful} \\
            }
        }
    }
    & & \\ \\

    \theutterance \stepcounter{utterance}  
    & & & \multicolumn{2}{p{0.3\linewidth}}{
        \cellcolor[rgb]{0.9,0.9,0.9}{
            \makecell[{{p{\linewidth}}}]{
                \texttt{\tiny{[GM$|$GM]}}
                \texttt{end game} \\
            }
        }
    }
    & & \\ \\

\end{supertabular}
}

\end{document}
